% ==============================================================================
% 1.4 MÉTRICAS (sobre tres software de referencia)
% ==============================================================================
% [PENDIENTE: DEFINIR CON EL EQUIPO] Los tres software de referencia. Cada uno lleva las
% mismas cinco subsecciones; las líneas de código se miden con cloc sobre el repositorio
% real (no se inventan cifras). Reemplazar «Software de referencia N» por el nombre elegido.
\subsection{1.4. Métricas}

% ---------------------------------------------------------------- Software 1
\subsubsection{1.4.1. Software de referencia 1}

\paragraph{1.4.1.1. Funcionalidades importantes.}

% [PENDIENTE]: Descripción general, repositorio, licencia y funcionalidades del software 1.

\paragraph{1.4.1.2. Captura de interfaces.}

% [PENDIENTE]: Capturas de pantalla del software 1.

\paragraph{1.4.1.3. Características del desarrollador.}

% [PENDIENTE]: Conteo de líneas de código (cloc) del software 1.

\paragraph{1.4.1.4. Métrica orientada al tamaño.}

% [PENDIENTE]: Tabla KLDC, costo, tiempo, esfuerzo, gente, calidad y productividad del software 1.

\paragraph{1.4.1.5. Aplicación de las métricas orientada a la función.}

% [PENDIENTE]: Puntos de función del software 1 (conteo, 14 factores Fi y PF final).

% ---------------------------------------------------------------- Software 2
\subsubsection{1.4.2. Software de referencia 2}

\paragraph{1.4.2.1. Funcionalidades importantes.}

% [PENDIENTE]: Ver 1.4.1.1.

\paragraph{1.4.2.2. Captura de interfaces.}

% [PENDIENTE]: Ver 1.4.1.2.

\paragraph{1.4.2.3. Características del desarrollador.}

% [PENDIENTE]: Ver 1.4.1.3.

\paragraph{1.4.2.4. Métrica orientada al tamaño.}

% [PENDIENTE]: Ver 1.4.1.4.

\paragraph{1.4.2.5. Aplicación de las métricas orientada a la función.}

% [PENDIENTE]: Ver 1.4.1.5.

% ---------------------------------------------------------------- Software 3
\subsubsection{1.4.3. Software de referencia 3}

\paragraph{1.4.3.1. Funcionalidades importantes.}

% [PENDIENTE]: Ver 1.4.1.1.

\paragraph{1.4.3.2. Captura de interfaces.}

% [PENDIENTE]: Ver 1.4.1.2.

\paragraph{1.4.3.3. Características del desarrollador.}

% [PENDIENTE]: Ver 1.4.1.3.

\paragraph{1.4.3.4. Métrica orientada al tamaño.}

% [PENDIENTE]: Ver 1.4.1.4.

\paragraph{1.4.3.5. Aplicación de las métricas orientada a la función.}

% [PENDIENTE]: Ver 1.4.1.5.
